\documentclass{article}
    
%Necessario per la traduzione di caratteri in encoding UTF-8 nel linguaggio di LaTeX
    \usepackage[utf8]{inputenc}
    
%Necessario per aggiungere determinati caratteri dell'alfabeto latino e altri simboli come quelli per le frecce
    \usepackage{lmodern,textcomp}

%Necessario per il font utilizzato
    \usepackage{tgheros}
    
%Necessario per la traduzione in italiano
    \usepackage[italian]{babel}

%Necessario per la definizione dei margini e della struttura della pagina
    \usepackage[nomarginpar, left=3cm,right=2cm,top=3cm,bottom=2cm]{geometry}

%Forza posizione immagini
\usepackage{float}

%Necessario per l'utilizzo dello stile di pagina "fancy" che permette l'header e il footer in ogni pagina e mostrare il numero di pagina sul totale
    \usepackage{fancyhdr}
    \usepackage{lastpage}

%Necessari per la creazione dei diagrammi e gli elementi disegnati della pagina
    \usepackage{tikz}
%Aiuta il pacchetto sopra, permettendo l'uso di nodi relativi alla pagina es(\node[anchor=south] at (current page text area.south))
    \usepackage{tikzpagenodes}
    
%XCOLOR
    \usepackage{xcolor}
    \newcommand{\blue}[1]{\textcolor{blue}{#1}}

%Necessario per permettere di utilizzare il parametro [H] per le tabelle e le figure in modo che vengano posizionate esattamente nel punto indicato 
    \usepackage{float}

%Istruzioni per indicare il numero di livelli di gerarchia per i quali generare la numerazione nel documento(Secnumdepth) e nella table of content (tocdepth)
% Valore 5 : Section -> SubSection -> SubSubSection -> Paragraph -> SubParagraph
    \usepackage{titlesec}
    \setcounter{secnumdepth}{5}
    \setcounter{tocdepth}{5}
%Istruzioni per specificare il desiderio di avere il rientro  e il grassetto anche per i paragrafi e sottoparagrafi
    \titleformat{\paragraph}
        {\normalfont\normalsize\bfseries}{\theparagraph}{1em}{}
    \titlespacing*{\paragraph}{0pt}{3.25ex plus 1ex minus .2ex}{1.5ex plus .2ex}
    \titleformat{\subparagraph}
        {\normalfont\normalsize\bfseries}{\thesubparagraph}{1em}{}
    \titlespacing*{\subparagraph}{0pt}{3.25ex plus 1ex minus .2ex}{.75ex plus .1ex}
    
%Necessario per poter includere hyperlink
    \usepackage[hidelinks]{hyperref}
%Tabelle
\usepackage{multirow}


\pagestyle{fancy}
    \fancyhf{}
    \lhead{A.C. Torre Archimede}
    \rhead{Relazione}
    \rfoot{Pagina \thepage / \pageref{LastPage}}
    
    
\title{
    \includegraphics[scale=0.5]{images/logo.png} \\
    \vspace*{1in}
    {
        \Huge \textbf{A.C. Torre Archimede}}\\
        \vspace*{0.25in}
        \textbf{Relazione Tecnologie Web}\\
        \vspace{0.2in}
        \textbf{Versione: 0.0.2}\\
        \vspace{0.1in}
    }

\author{
    \begin{tabular}[t]{c@{}c}
        \hline
        \\
        Michele Bettin \\
        Marco Marchiante \\
        Stefano Meneguzzo\\
        Edoardo Retis\\
    \end{tabular}
    \vspace*{0.5in} \\
    Dipartimento di Matematica \\
    \textbf{Università degli Studi di Padova} \\
} 
\date{\today}


%--------------------Make usable space all of page
\setlength{\oddsidemargin}{0in} \setlength{\evensidemargin}{0in}
\setlength{\topmargin}{0in}     \setlength{\headsep}{.25in}
\setlength{\textwidth}{6.5in}   \setlength{\textheight}{8.5in}
%--------------------Indention
\setlength{\parindent}{1cm}

\begin{document}
{
    %\fontfamily{qhv}
    \selectfont
    \maketitle
    \newpage
    \tableofcontents
    \newpage

    \section{Abstract}
\subsection{Introduzione}
Il progetto illustrato da questa relazione ha lo scopo di creare un sito internet che rappresenta una fittizia società calcistica, denominata "Associazione Calcistica Torre Archimede". Tale società ha sede in Italia, più precisamente a Padova e opera sotto la giurisdizione della Federazione Italiana Giuoco Calcio (FIGC). La società ha una squadra di calcio che partecipa a campionati italiani organizzati dalla FIGC, uno stadio, denominato "TorreArchimede Stadium", sede della stessa e che usa per allenare la propria squadra, per disputare gare "in casa", e una tifoseria che ogni settimana si riunisce al botteghino per acquistare un biglietto.\\
L'obiettivo del sito è di essere un punto di riferimento sia per la società calcistica sia per i fan che navigano su Internet, cercando di raggiungere il maggior numero di persone. Infatti, il sito è usato per rendere partecipi dell'attività della società ai tifosi che già frequentano lo stadio e, al tempo stesso, permettere ai nuovi tifosi di venire a conoscenza del mondo calcistico del Torre Archimede.\\
Il sito fornisce degli strumenti per permettere a tutte le persone interessate di acquistare online i biglietti di ingresso al campo da gioco, alleggerendo il lavoro di vendita ai tradizionali botteghini.




    \newpage
    \section{Analisi}
    \subsection{Analisi dell'Utenza}
    Il target principale del sito è la tifoseria dell'AC Torre Archimede che, pur essendo abituata al mondo del calcio e familiare con lo stadio della squadra, potrebbe non essere pratica di strumenti informatici. Per questo motivo, il sito internet deve tenerne conto, rendendo il contenuto il più chiaro e semplice possibile.\\
    Un altro punto da considerare sull'utenza del sito è la sua provenienza: la squadra dell'AC Torre Archimede partecipa soltanto a campionati minori, senza mai operare a livello internazionale, presumendo che l'utenza sia italiana, per la maggior parte di loro, residente in zone limitrofe allo stadio. Per questo motivo, la lingua usata nel sito deve essere quella italiana (ad eccezione dei comuni termini di internet, ad esempio, Home, News, Login, ecc.), non richiedendo più versioni alternative in diverse lingue.

    \subsection{Base Informativa}
    Secondo un analisi preliminare del tipo di utenti, le informazioni che un utente si aspetta di trovare e le domande che il sito deve rispondere quando esso vi si accede riguardano l'attività sportiva della società, in particolar modo:
    
    \begin{itemize}
        \item Lo stadio (Dove si trova? Come è fatto?);
        \item La squadra (Chi sono i calciatori? e i dirigenti? A quali eventi partecipano?);
        \item Le gare (Quando giocano? In che stadio? Contro chi? Quanto costa un biglietto?).
    \end{itemize}
    
    \noindent
    Tenendo conto di queste considerazioni, per soddisfare il fabbisogno informativo dell'utente, il sito deve proporre informazioni che riguardano gli argomenti sopra citati, utilizzando come strumenti comunicativi:
    
    \begin{itemize}
        \item Testo descrittivo, se si vuole utilizzare descrizioni per approfondire determinati argomenti (ad esempio, descrivendo come è fatto lo stadio della squadra);
        \item Articoli, se si vuole tenere aggiornato l'utente su fatti che accadono nella sfera quotidiana della società;
        \item Elenchi e tabelle, se si vuole comunicare all'utente una raccolta di dati riferiti ad un certo argomento.
    \end{itemize}


    \subsection{Attori}
    Basandosi sulle necessità del sito, tutti gli utenti si possono raggruppare in tre tipologie: gli utenti non autenticati, gli utenti autenticati e gli amministratori.
    
        \subsubsection{Attore utente non autenticato}
        
        Gli utenti non autenticati sono persone che possono navigare sul sito, ma hanno qualche limitazione. Possono visitare la Homepage, guardare le notizie pubblicate dalla società, leggere la descrizione sullo stadio, consultare l'elenco delle gare in programma, ma non possono acquistare biglietti oppure accedere a sezioni apposite per altri tipi di utenti. Questo tipo di utente può registrarsi al sito per ottenere delle credenziali d'accesso, utili per poter fare l'autenticazione, oppure autenticarsi se ne è già in possesso.
        
        \subsubsection{Attore utente autenticato}
        
        Gli utenti autenticati sono persone che hanno inserito le proprie credenziali d'accesso e si sono autenticati correttamente. Oltre a poter navigare in tutte le aree del sito, possono acquistare biglietti, scegliendo la gara e il settore dove guardarla, e accedere ad una sezione personale dove vengono visualizzati i dati anagrafici dell'utente e l'elenco dei suoi biglietti acquistati.\\ L'utente autenticato può anche togliere l'autenticazione per ritornare ad essere un utente non autenticato.
        
        \subsubsection{Attore amministratore}
        
        Gli amministratori sono utenti autenticati al sito con delle credenziali d'accesso apposite che gli permettono l'ingresso ad una sezione riservata esclusivamente a loro. In questa area, hanno strumenti per:
        \begin{itemize}
            \item Aggiungere articoli inerenti alla società;
            \item Aggiungere o togliere le gare del calendario.
        \end{itemize}
        Anche loro hanno la possibilità di navigare nel sito e fare acquisti di biglietti.
        

    \subsection{Requisiti funzionali}
    Il primo requisito funzionale che il nostro sito deve adempiere è avere una comunicazione efficace delle informazioni da esporre; cioè organizzare il contenuto in modo tale da poter essere facilmente compreso dalla maggior parte degli utenti che navigano in Internet. Per soddisfare questo requisito, si è sviluppato il sito concentrandosi di più sul contenuto da esporre, sviluppando un design minimale, ma efficace e adattabile a qualsiasi dispositivo.\\
    Un altro requisito da soddisfare è mettere a disposizione dell'utente una piattaforma che gli permetta di fare acquisti online per comprare dei biglietti; a questo scopo si è implementato all'interno del sito una serie di strumenti per:
    
    \begin{itemize}
        \item Controllare le partite future e/o da acquistare;
        \item Acquistare un biglietto;
        \item Controllare i biglietti acquistati;
        \item Cancellare e chiedere eventualmente il rimborso per un biglietto acquistato.
    \end{itemize}
    
    \noindent
    L'ultimo requisito da soddisfare, ma non meno importante, è avere un sito facilmente navigabile e accessibile a tutti. Il sito deve essere accessibile per evitare di limitare la navigazione alle persone con difficoltà nel distinguere i colori o alle persone non vedenti.\\
    Per soddisfare questo requisito sono state studiate delle strategie per permettere sia a persone normodotate che diversamente abili di riuscire a trovare facilmente tutte le informazioni e i contenuti all'interno del sito.
    
    


    \newpage
    \section{Organizzazione interna}
    \subsection{Divisione dei compiti}
    Tutti i membri del gruppo hanno avuto modo di toccare tutte le tipologie di lavoro, alternandosi a vicenda nello svolgere i diversi compiti.\\
    Nell'elenco sottostante riportiamo le parti più sostanziali svolte da ogni componente:
    \begin{itemize}
        \item \textbf{Marco}: Back-end PHP, CSS Desktop e Javascript;
        \item \textbf{Michele}: Database, Back-end PHP e CSS Desktop;
        \item \textbf{Edoardo}: Struttura pagine HTML e Relazione;
        \item \textbf{Stefano}: Javascript e CSS Mobile/Tablet/Stampa.
    \end{itemize}
    
    \subsection{Tempistiche}
    Lo sviluppo del progetto è inizialmente cominciato a Dicembre 2021, ma è stato deciso di ricominciarlo da zero a Febbraio 2022, in modo da raggiungere una qualità che potesse soddisfare tutti i membri del gruppo e rispettare tutti i requisiti richiesti dalla consegna.\\
    Da Febbraio fino a Giugno le scadenze sono state rispettate con continuità. Ci sono stati solo dei ritardi nelle ultime settimane, a causa della mancanza di tempo del gruppo, concentrato a studiare per la sessione estiva.
    In conclusione, ci riteniamo soddisfatti del rispetto delle scadenze.

    \newpage
    \section{Tecnologie}
    \subsection{Ambiente di sviluppo}
    Lo sviluppo è stato effettuato su Windows 10/11 e in parte su Ubuntu 21.10 da parte di 4 diversi sviluppatori e con l’utilizzo di Visual Studio Code come IDE.
    Per questa ragione la maggior parte dello sviluppo, del backend è stato eseguito in un ambiente virtuale (un container docker in particolare) con una configurazione il più possibile simile all’ambiente di deploy.
    Ciò unito a uno sviluppo collaborativo tramite Github ha avuto come conseguenza una bassa  incompatibilità con l’ambiente di deploy ee il vantaggio di non dover costantemente fare test di caricamento sul server di laboratorio.
    \subsection{Tecnologie utilizzate}
    Le tecnologie utilizzate per la realizzazione del sito sono state in parte influenzate dall’ambiente finale di deploy (quale il server di laboratorio), in particolare gli applicativi lato Server: PHP (v. 7.4.3), Apache2 (v. 2.4.41), MariaDB (v.10.3.34) e Ubuntu (20.02.4 LTS).\\
    A lato Client è stato utilizzato ovviamente JavaScript e per quanto riguarda HTML e CSS si è optato per HTML5 e CSS3. Nella scelta tra XHTML 4.1 e HTML5 e tra CSS 2 e 3 si è optato per l’ultima versione poichè pur essendo tecnologie nuove e non ancora pienamente supportate, varie fonti (W4C survey e \blue{\url{caniuse.com}}) riportano una diffusione tale da non creare problemi di compatibilità anche con i browser più vecchi a, patto di non utilizzare componenti troppo avanzate dello standard. Oltre alle informazioni reperite in rete abbiamo eseguito anche i nostri piccoli test che hanno confermato le nostre ricerche.\\
    Una volta fatti i conti con la realtà, è risultato che anche i browser più vecchi supportano discretamente lo standard HTML5 (fatta eccezione per IE 6 e 9, che a volte è sono un po’ restrittivi ma in casi non troppo complicati è sempre possibile un workaround), mentre i problemi sono più forti con CSS3 che è ben supportato dalla maggior parte dei browser attualmente in uso, ma ha grosse difficoltà nei browser più vecchi come IE e Opera Mini che supportano molto poco le caratteristiche principali del nuovo standard come flex, grid e attributi avanzati.\\
    Fatte queste considerazioni, è stato deciso di usufruire dei vantaggi di pulizia e raffinatezza offerti dal nuovo standard CSS3.\\
    Per quanto riguarda le versioni del software client (quindi JavaScript e browser), il sito è stato prodotto sulla latest release di Google Chrome ma testato frequentemente su altre versioni spaziando inoltre tra diversi browser (vedi capitolo dei \blue{\hyperref[Test]{Test di accessibilità}}).
    \subsection{Configurazione e installazione}
    Il sito è già configurato e per eseguirlo è sufficiente posizionare i files all'interno della root directory del web server.\\
    Per configurare il database, è sufficiente configurare nel file \texttt{/php/actions/DBAccess.php} i parametri di configurazione e importare il file SQL nel DBMS.\\
    Bisogna poi modificare adeguatamente i files \texttt{.htaccess} e \texttt{/php/helpers/Session.php} (righe 28 e 46), riportando il path assoluto rispetto alla web root del file \texttt{/php/error.php}, per consentire di visualizzare gli errori a prescindere da dalla pagina corrente.
    \newpage
    \section{Progettazione}
    \subsection{Schemi di ordinamento}
    \subsection{Emotional Design}
    \subsection{Accessibilità}
    \subsection{Usabilità}


    \newpage
    \section{Base di Dati}
    \subsection{Scelta delle tabelle}
    \subsection{Identificatori primari}
    \subsection{Forma normale}
    \subsection{Utilizzo atteso}

    
    \newpage
    \section{Struttura}
    \subsection{Separazione tra struttura, comportamento e presentazione}
    Uno dei problemi che si è dovuto risolvere per la scrittura del codice HTML è stato quello di evitare l'uso di codice "inline" di altri linguaggi di programmazione all'interno dei file \texttt{.html}, nell'ottica di  garantire il principio di separazione tra struttura, presentazione e comportamento. L'idea è stata quella di frammentare il codice in più parti ricostruendolo successivamente con delle chiamate dinamiche, ogni qualvolta venga caricata una pagina, utilizzando appositi file di codice HTML denominati \textit{presets}; alcuni \textit{presets} contengono il mero contenuto da stampare a video, altri invece, chiamati \textit{template}, forniscono del codice utile a costruire la struttura di una pagina web. L'implementazione di questa idea è supportata anche da stringhe racchiuse da una coppia di percentuali, chiamate \textit{placeholders}, necessarie per inserire dinamicamente dei dati generati dal server durante l'invocazione di una pagina (vedi capitolo \blue{\hyperref[Comportamento]{Comportamento}}).\\ I vantaggi che porta questo approccio sono:
    \begin{itemize}
        \item Migliore manutenibilità del codice;
        \item Codice più pulito.
    \end{itemize}
    
    \subsection{Organizzazione generale del codice}
    Per l'implementazione del codice HTML, si è deciso di seguire le seguenti regole:
    \begin{itemize}
        \item Per creare una struttura più semantica ed evitare tag \texttt{<div>} incorretti, incapsulati uno dentro l'altro, è stato utilizzato il tag \texttt{<section>};
        \item Per incapsulare porzioni di codice che rappresentano aiuti per la navigazione è stato utilizzato il tag \texttt{<nav>};
        \item Le varie opzioni per gli aiuti alla navigazione sono state organizzate con i tag \texttt{<li>} e \texttt{<ul>};
        \item Per rappresentare le intestazioni principali di pagina, sono stati utilizzati i tag \texttt{<h1>} e \texttt{<h2>}; per rappresentare le intestazioni di contenuto si è utilizzato il tag \texttt{<h3>};
        \item Viene utilizzato il tag \texttt{<abbr>} per rappresentare abbreviazioni;
        \item Per indicare termini in inglese è stato utilizzato l'attributo \texttt{lang=”en”}.
    \end{itemize}
    
    \subsection{Altre regole strutturali}
    Di seguito, sono elencate alcune regole strutturali specifiche per determinate pagine:
    \begin{itemize}
        \item Nella pagina della \textit{Home} è presente una sezione che rappresenta un'anteprima dei più recenti articoli presenti nella pagina \textit{News}; tale blocco è stato incapsulato dentro un tag \texttt{<aside>}. Questa decisione è stata presa perché il contenuto all'interno del tag non è indispensabile per la semantica della \textit{Home} e serve solo a supportare la navigazione nel sito;
        \item Nella pagina delle \textit{News}, ma anche nella pagina \textit{Home} tutti gli articoli sono racchiusi dentro un tag \texttt{<article>} per identificare porzioni di testo che sono indipendenti tra di loro;
        \item Per organizzare la struttura di ogni tabella, l'intestazione è stata incapsulata dentro il tag \texttt{<thead>}, mentre il corpo è stato incapsulato dentro il tag \texttt{<tbody>};
        \item Nella pagina \textit{Registrazione} è stato usato il tag \texttt{<fieldset>} per suddividere la form in due sezioni: la sezione per compilare i dati anagrafici dell'utente e la sezione per compilare i dati di accesso all'area personale del sito;
        \item Nella pagina dedicata all'utente, la struttura per visualizzare le informazioni anagrafiche è stata costruita usando elenchi di definizioni, usando i tag \texttt{<dl>}, \texttt{<dt>} e \texttt{<dd>};
        \item Se un utente è autenticato come amministratore nell'area personale, è presente un pulsante per passare alla dashboard di amministrazione, che contiene: le form per inserire gli articoli, un pulsante per inserire una nuova partita e, per ogni partita già presente, la possibilità di modificarla o eliminarla.
    \end{itemize}
    \newpage
    \section{Presentazione}
\label{Presentazione}
    \subsection{Suddivisione dei file}
    Per gestire gli aspetti visivi del sito internet, sono stati prodotti quattro file:
    \begin{itemize}
        \item \texttt{style.css} per gli schermi desktop o con pixel maggiori di 1061px;
        \item \texttt{tablet.css} per i dispositivi tablet o per schermi con pixel minori di 1061px;
        \item \texttt{mobile.css} per dispositivi smartphone o per schermi con pixel minori di 810px;
        \item \texttt{print.css} per stampe con stampanti tradizionali.
    \end{itemize}
    La costruzione di tutti questi file \texttt{.css} hanno lo scopo di ottenere un design capace di adattarsi a qualsiasi dispositivo. A supporto di questa implementazione le misure usate per scrivere le regole CSS sono misure relative come \texttt{em} o percentuali, in modo che l'intero sito possa scalare adeguatamente adattandosi alle scelte dell'utente.
    
    \subsection{CSS Desktop}
    Le caratteristiche del design per gli schermi grandi sono elencati di seguito:
    \begin{itemize}
        \item Il titolo, il logo e la barra di navigazione si espandono in orizzontale, occupando tutto lo schermo; tutti i link della barra di navigazione sono allineati in orizzontale;
        \item Nella \textit{Home} la sezione relativa all'anteprima degli articoli è spostata a destra;
        \item Nella pagina dedicata alle \textit{News} tutti gli articoli sono raggruppati con una grid per evitare troppo scrolling verticale;
        \item L'intestazione e i record delle tabelle si espandono in orizzontale occupando tutto lo schermo a disposizione.
    \end{itemize}
    
    \subsection{CSS Tablet e mobile}
    Le caratteristiche del design per gli schermi di dispositivi portatili sono elencati di seguito:
    \begin{itemize}
        \item Il titolo del sito si toglie lasciando spazio alla visualizzazione soltanto del logo che funge da titolo del sito;
        \item La barra di navigazione viene ridimensionata allineando in verticale i link di navigazione nel sito con i link per l'account dell'utente; se lo schermo si riduce ulteriormente la barra di navigazione si toglie lasciando spazio a due pulsanti che, se premuti, fanno scendere un menù a tendina che mostra i link della barra;
        \item Nella \textit{Home} la sezione relativa all'anteprima degli articoli è spostata al centro seguendo il normale flusso del contenuto;
        \item Nella pagina dedicata alle \textit{News} tutte gli articoli sono visualizzati in verticale uno sotto l'altro;
        \item Il design delle tabelle è stato ridimensionato per trasformarsi in maniera elegante ed evitare lo scrolling orizzontale. Ogni record della tabella viene visualizzato come un blocco dove ogni cella ha la propria intestazione e il proprio dato informativo; le intestazioni incapsulate dentro il tag HTML \texttt{<thead>} si tolgono e ciò impatta sullo scroll verticale aumentando la lunghezza della pagina da scorrere, ma non crea frustrazione perché è accettato dall'utente;
        \item Le form restano invariate.
    \end{itemize}
    
    \subsection{CSS per stampe}
    Le caratteristiche del design per le stampe sono elencati di seguito:
    \begin{itemize}
        \item La barra di navigazione viene nascosta lasciando spazio alla visualizzazione del logo, del titolo e della breadcrumb;
        \item Tutti gli elementi grafici sono stati nascosti lasciando soltanto il mero contenuto con un font con grazie adatto alle stampe;
        \item Le tabelle sono visualizzate normalmente in orizzontale con le celle delimitate da bordi neri e al loro interno il contenuto.
    \end{itemize}
    \newpage
    \section{Comportamento}
\label{Comportamento}
    \subsection{Sviluppo del codice}
    Lo sviluppo dell’aspetto comportamentale del sito è avvenuto in PHP (lato Server) e JavaScript (lato Client).
    Ogni qualvolta il server web registra una richiesta per uno specifico documento, viene eseguito uno script PHP che genera la pagina.
    \subsection{Comportamento lato Server}
    In generale ogni pagina è composta da alcuni elementi fissi come intestazione, barra di navigazione e footer che non cambiano mai e un corpo che differisce da pagina a pagina e molte volte differisce anche all’interno della stessa pagina a seconda della situazione; quindi PHP si occupa di assemblare queste componenti in un documento HTML valido e modificare il corpo della pagina a seconda di ciò che le circostanze richiedono.
    Ad esempio l’area personale dell’utente dovrà mostrare le informazioni relative all’utente attualmente autenticato sulla piattaforma oppure la pagina \textit{Home} dovrà mostrare l’anteprima dei prossime eventi e gli ultimi articoli. Nessuna di queste informazioni e fissa e le pagine vi si devono adattare.
    L’ultimo compito del server è di gestire gli errori come ad esempio tentativi di accesso a risorse ad accesso ristretto (risorse amministrative, file di backend, risorse per utenti autenticati etc…) o problemi a lato server (un malfunzionamento del server web o del database), condizione per cui sono previsti più casi:
    \begin{itemize}
        \item In caso di deliberato accesso (il sito non lo permette quindi tramite URL) a risorse ad accesso limitato o in caso di malfunzionamento del database (fallimento di connessione o errore di query) avviene una ridirezione ad una pagina di errore che fornisce gli strumenti per continuare la navigazione ed un messaggio di errore esplicativo;
        \item In caso di tentato accesso ad una funzionalità non applicabile (come \textit{Login}/\textit{Registrazione} per un utente autenticato o \textit{Logout} per utente ospite) avviene la normale ridirezione all’\textit{Area personale}/\textit{Home} senza che avvenga alcuna comunicazione all’utente.
    \end{itemize}
    L’interattività di una pagina in senso “PHP” è limitata a quelle interazioni che richiedono in qualche modo un cambio o “refresh” della pagina corrente e per qualunque altra interazione è necessario ricorrere a JavaScript poiché operando a livello client è in grado di interagire con il DOM e di conseguenza interagire con la pagina in maniera più "soft" di PHP.
    \subsection{Comportamento lato Client}
    Nel nostro progetto JavaScript ha utilità marginale, nel senso che la maggior parte dei comportamenti interattivi è gestita come descritto precedentemente, ma non per questo è meno importante. Infatti viene utilizzato per modificare dinamicamente il totale durante il processo di checkout e, molto più importante, si occupa di validare l’input nelle form e fornire un feedback all’utente circa l’esito dell’operazione.\\
    Il primo comportamento risiede in un file a se stante e raccogliendo l’evento \texttt{onchange} sul numero di articoli ordinati è in grado di interagire anche con un apposito paragrafo e la lista dei prezzi; in fine i rispettivi dati vengono elaborati e presentati sotto forma di un totale fittizio che viene ricalcolato dal server in fase di ultimazione del checkout.
    Il secondo è più complesso poiché la validazione delle form viene effettuata in più momenti e tutto deve integrarsi con gli errori restituiti dal server aumentando la complessità dell’intero sistema.\\
    La prima validazione viene effettuata “in tempo reale” con lo scopo di informare tempestivamente l’utente di un possibile errore nella compilazione: ogni volta che un campo viene compilato, viene anche validato tramite una specifica funzione JavaScript, attivata dell’evento \texttt{onfocusout}, e si occupa di valutare la correttezza dell’input secondo parametri specifici, generare l’errore e stamparlo a schermo in maniera accessibile per le tecnologie assistive e correttamente presentato per essere facilmente identificato nella pagina.\\
    La seconda validazione viene effettuata prima che la form venga inviata come controllo di sanità generale poiché è possibile che l’utente non abbia compilato tutti i campi e rimbalzare il controlllo al server è poco usabile per l’utente: quando una form viene inviata al server, l’evento \texttt{onsubmit} della form aziona una particolare funzione che effettua i controlli per tutta la form e consente l’invio solo se nessun errore si è verificato. Per evitare la duplicazione del codice questa classe di funzioni  di validazione si limita ad eseguire le stesse funzioni del punto precedente, che aggiornano a loro volta una variabile per registrare la presenza di errori nella compilazione della form. Questa variabile viene poi letta da questa funzione (e salvata), resettata e in base al suo valore la form verrà inviata o meno.
    \newpage
    \input{10_Accessibilità}
    \newpage
    \input{11_Usabilità}
    \newpage
    \section{Conclusioni}
Il progetto è stato portato a termine entro i tempi previsti senza incontrare particolari problemi ed ogni membro del gruppo ha saputo gestire bene le tecnologie utilizzate per creare un sito internet efficiente dal punto di vista dei contenuti, dell'accessibilità e dell'usabilità. Sebbene sia stato usato HTML5 e CSS3 alcune funzioni di questi linguaggi, come ad esempio le canvas, non sono state utilizzate. Questa scelta è stata fatta perché l'utenza e le richieste del sito necessitano di poche funzionalità.\\
Infine sono stati eseguiti sul sito tutti i test di accessibilità riportati dentro \blue{\url{https://web.math.unipd.it/accessibility-dev/}}.
    \newpage
}

\end{document}